\documentclass{article}
\usepackage{graphicx}

\title{AutoCTM Software Specification}
\author{Oscar Ashburn}
\date{April 2026}

\begin{document}

\maketitle

\begin{abstract}
This document specifies the functional and technical requirements for AutoCTM 
(Automated Certificate Transparency Monitor), a tool written in Go for 
monitoring Certificate Transparency (CT) logs. It defines the system's 
expected behaviour, interfaces, and constraints, serving as the authoritative 
reference for the design and implementation of AutoCTM.
\end{abstract}

\newpage

\section{Introduction}

\subsection{What is Certificate Transparency}

Certificate Transparency (CT) is a system designed to make TLS/SSL certificate 
issuance public, auditable, and verifiable. It helps prevent mis-issued or 
rogue certificates and increases trust in HTTPS.

\subsection{The Problem}

Certificate Authorities (CAs) issue TLS/SSL certificates for domains. 
Sometimes, a CA can mistakenly issue a certificate to the wrong party, or a 
malicious actor could obtain one. Without transparency, domain owners might 
not know about these certificates until it is too late.

\begin{itemize}
    \item CAs can mistakenly or maliciously issue certificates for domains 
    they should not.
    \item Domain owners may have no visibility into certificates issued in 
    their name.
    \item The window of unawareness is the dangerous part. An attacker could 
    use a fraudulent certificate before anyone notices.
\end{itemize}

\subsection{How CT Works}

Certificate Transparency solves the problem of rogue certificate issuance by 
requiring every CA to submit certificates to public, append-only logs before 
browsers will trust them. This means anyone can monitor these logs to spot 
unauthorised certificates for their domain, and browsers will reject any 
certificate that lacks proof of being logged.

\section{System Architecture Overview}

AutoCTM is implemented as three cooperating processes that communicate through 
a local control socket and a shared database.

\begin{itemize}
    \item \textbf{CLI} (\texttt{autoctm}) — The operator-facing interface used 
    to control the system.
    \item \textbf{Broker} (\texttt{autoctm-broker}) — A persistent background 
    process responsible for orchestration and instance management.
    \item \textbf{Monitor Instance} (\texttt{autoctm-instance}) — A worker 
    process responsible for monitoring CT logs and processing certificate 
    entries.
\end{itemize}

The broker acts as the control plane of the system, while monitor instances 
perform the actual log monitoring work. Instances interact directly with the 
database and continue monitoring even if the broker becomes unavailable.

\section{Design Description}

AutoCTM is written in Go and is composed of three decoupled processes: a 
persistent broker daemon, one or more monitor instances, and a lightweight 
operator CLI.

The system must support the following functional requirements.

\subsection{Broker Daemon}

AutoCTM shall run a persistent broker daemon that acts as the backbone of 
the system. The broker shall:

\begin{itemize}

\item Own and maintain a Unix domain socket at a well-known path used for 
all CLI-to-broker and instance-to-broker communication.

\item Maintain an instance registry tracking each instance's ID, PID, and 
current status.

\item Persist the instance registry to the database so that it survives 
broker restarts.

\item On restart, reload the instance registry from the database and wait 
for running instances to reconnect and re-register.

\item Spawn new monitor instance processes on operator request.

\item Route commands from the CLI to the appropriate instance.

\item Start automatically when the CLI is launched and no broker is found.

\end{itemize}

\subsection{CLI Interface}

AutoCTM shall provide a command-line interface for operator control. The CLI 
shall:

\begin{itemize}

\item On launch, connect to the broker socket. If no broker is found, 
automatically spawn one before presenting the command prompt.

\item Operate in one of two contexts: \textit{global context} (no instance 
selected) and \textit{instance context} (an active instance is selected via 
\texttt{/set-instance}).

\item Support the following commands in global context:

\begin{itemize}

\item \texttt{/start} --- Spawn a new monitor instance.

\item \texttt{/get-instances} --- List all instances and their status.

\item \texttt{/set-instance <instanceID>} --- Set the active instance. 
Subsequent commands are routed to this instance.

\item \texttt{/add <domain>} --- Add a domain to the shared watchlist.

\item \texttt{/remove <domain>} --- Remove a domain from the shared watchlist.

\item \texttt{/shutdown} --- Gracefully shut down all instances and the broker.

\end{itemize}

\item Support the following commands in instance context:

\begin{itemize}

\item \texttt{/add-log <url>} --- Add a CT log to the instance configuration.

\item \texttt{/remove-log <url>} --- Remove a CT log from the instance.

\item \texttt{/pause} --- Pause the instance's polling loop.

\item \texttt{/resume} --- Resume a paused instance.

\item \texttt{/status} --- Display the instance's current status.

\item \texttt{/shutdown} --- Shut down the active instance only.

\item \texttt{/clear-instance} --- Return to global context.

\end{itemize}

\item May be closed at any time. Running instances and the broker shall be 
unaffected. On relaunch the CLI shall reconnect to the existing broker.

\end{itemize}

\subsection{Monitor Instances}

AutoCTM shall support spawning one or more independent monitor instances as 
background processes. Each instance shall:

\begin{itemize}

\item Load its configuration from the shared database using its instance ID.

\item Monitor one or more RFC 9162-compliant CT logs.

\item Poll logs for new Signed Tree Heads (STHs) on a configurable interval.

\item Verify STH signatures before processing entries.

\item Fetch new entries since the last known STH.

\item Verify Merkle consistency proofs between successive STHs.

\item Support monitoring of multiple logs concurrently using independent 
goroutines.

\item Continue monitoring even if the broker becomes unavailable.

\end{itemize}

\subsection{Certificate Storage}

AutoCTM shall persist all observed certificate entries to a shared PostgreSQL 
database. The system shall:

\begin{itemize}

\item Store the leaf hash, log ID, timestamp, subject, SANs, issuer, serial 
number, and expiry for each observed certificate.

\item Enforce uniqueness of certificates based on the leaf hash.

\item Support querying stored certificates by domain, issuer, and date range.

\end{itemize}

\subsection{Certificate Revocation Checking}

AutoCTM shall check observed certificates against Certificate Revocation 
Lists (CRLs). The system shall:

\begin{itemize}

\item Fetch and cache CRLs for issuers relevant to observed certificates.

\item Refresh CRLs on a configurable interval.

\item Flag and record any certificate whose serial number appears in a 
current CRL.

\end{itemize}

\subsection{Domain Watchlist Alerting}

AutoCTM shall allow operators to configure a shared watchlist of domains to 
monitor for unexpected certificate issuance. The system shall:

\begin{itemize}

\item Match observed certificates against the watchlist by exact domain and 
wildcard SAN.

\item Raise an alert for any certificate covering a watched domain.

\item Allow watchlist entries to be added and removed via the CLI without 
restarting any instance or the broker.

\end{itemize}

\subsection{Log Integrity Auditing}

AutoCTM shall detect and report misbehaviour in monitored logs as defined 
in RFC 9162. The system shall:

\begin{itemize}

\item Detect when entries are not incorporated into the log within the 
declared Maximum Merge Delay (MMD).

\item Report entries that remain unavailable for an extended period after 
appearing in an STH.

\item Record and report any consistency proof verification failure.

\item Pause polling of a log on consistency proof failure pending operator 
review.

\end{itemize}

\subsection{Alert Notification}

AutoCTM shall deliver alerts to operators via configurable notification 
channels. The system shall:

\begin{itemize}

\item Support at least one outbound notification channel (e.g. email or 
webhook).

\item Include in each alert the affected certificate details, the triggering 
log, and the reason for the alert.

\item Avoid duplicate alerts for the same certificate and alert type within 
the same log.

\end{itemize}

\section{Non-Functional Requirements}

\subsection{Performance}

The system shall support monitoring multiple CT logs concurrently without 
blocking other log polling operations.

\subsection{Reliability}

Monitor instances shall continue polling CT logs and writing results to the 
database even if the broker process becomes unavailable.

\subsection{Scalability}

The system shall support running multiple monitor instances simultaneously to 
distribute monitoring load.

\subsection{Security}

Database credentials and sensitive configuration shall not be exposed through 
the CLI or stored in plaintext within the application source code.

\end{document}